% s-319-p319-m003; original PDF 319, 320
\Aphorism{s-319-p319-m003}{376}
Wie würde auf moderne Europäer die Heraufkunft metempsychotischer Vorstellungen wirken? Eine Frage, die einmal acut werden könnte, gesetzt, dass der Buddhismus, wie ihm von Einigen geweissagt wird, nicht nur die Salons, sondern auch die Kirchen und Schulen des erlöschenden Christenthums erbte. Man darf in diesem Falle einen neuen Beleg zu der alten Wahrheit erwarten, dass jede Mitleidsmoral als Multiplicator des Elends wirkt und zu dem vorhandenen physischen Leiden ein psychisches fügt, dessen Gefährlichkeit für Leben und Lebensenergie ins Uncontrollirbare wächst. Zunächst einmal würde das absichtlich verursachte Leiden, dessen relativer Betrag im Verhältniss zum Gesammtquantum des Leidens übrigens gering ist, durch eine Mitleidsreligion mit Seelenwanderungsglauben nur unwesentlich abnehmen. Der neurasthenische Feigling, der mit schlotternden Knieen Thiere oder Kinder foltert, würde vielleicht durch die Erkenntniss, sich ins eigene Fleisch zu schneiden, abgeschreckt; der Imperator, der Tausende aufs Schlachtfeld führt, würde das Bewusstsein, das Schicksal dieser Tausende persönlich durch Wiedergeburt zu erleiden, nur als eine willkommene Steigerung in seinen Machtrausch hineinnehmen. Das unvermeidliche Leiden dagegen, als Krankheit, Tod, Kampf ums Dasein in Thier- und Menschenwelt, träte fortan mit jenem psychischen Multiplicator auf, der die eigentlich lebensfeindliche, pessimistische Interpretation des Leidens zum Wahnsinn steigert. Wer irgendwie zum Indifferentismus neigt, fände sich in einer Lehre gerechtfertigt, die die Grenzen des Individuums aufhebt und alles Erduldete und Erduldbare zum gemeinsamen Loose Aller macht; der europäische décadent, der ohnehin alle erdenklichen Affinitäten zum Nihilismus hat, würde ja keinen Finger mehr rühren, um „dem Bösen zu widerstehen“, wenn er auf das Meer universalen Leidens blickt, das er mit oder ohne individuelle Verschuldung, so oder so doch austrinken muss. Vegetatives Starren und Dämmern, gelegentlich asketische Versuche, sich auf die höchsten Leidensgrade gewissermassen einzuüben: lockt uns diese indische Perspective, uns Europäer?
\par

